\documentclass[a4paper]{amsart}

\usepackage{color}
\usepackage[colorlinks=true, citecolor=red]{hyperref}  %needs to come before amsrefs package to work with \cite{}*{}
\usepackage{amssymb,latexsym}
\usepackage{amscd}  %commutative diagrams
\usepackage[alphabetic]{amsrefs}
%\usepackage[mathscr]{eucal}
\usepackage{pdfsync}
\usepackage{graphicx}
\usepackage{epstopdf}
\DeclareGraphicsRule{.tif}{png}{.png}{`convert #1 `dirname #1`/`basename #1 .tif`.png}
\usepackage[all]{xy}
\usepackage{titletoc}
%\usepackage{txfonts}   %  \fint  for integral average

%\usepackage{showkeys}  %options notref  notcite 

%\input{rgb.tex}

%\textcolor{DarkRed}{text in DarkRed}


\CompileMatrices

\theoremstyle{plain}
\newtheorem{theorem}{Theorem}[section]
\newtheorem{lemma}[theorem]{Lemma}
\newtheorem{proposition}[theorem]{Proposition}
\newtheorem{corollary}[theorem]{Corollary}

\theoremstyle{definition}
\newtheorem{definition}[theorem]{Definition}%[section]


\theoremstyle{remark}
\newtheorem{defn}[theorem]{Definition} 
\newtheorem{example}[theorem]{Example}
\newtheorem{remark}[theorem]{Remark}
\newtheorem{discussion}[theorem]{Discussion}

%\numberwithin{equation}{section}

%\newcommand{\abs}[1]{\lvert#1\rvert}

\DeclareMathOperator{\Lip}{Lip}
\DeclareMathOperator{\dist}{dist}
\DeclareMathOperator{\expected}{\mathbb{E}}
\DeclareMathOperator{\prb}{Prob}
\DeclareMathOperator{\spt}{spt}
\DeclareMathOperator{\lap}{\Delta}
\DeclareMathOperator{\Tr}{Tr}
\DeclareMathOperator{\diam}{diam}

\raggedbottom

\addtolength{\textwidth}{2cm}
\addtolength{\oddsidemargin}{-1cm}
\addtolength{\evensidemargin}{-1cm} 
\addtolength{\topmargin}{-1cm} 
\addtolength{\textheight}{2cm}

 
%%% BEGIN DOCUMENT
\begin{document}

\title{Perron Frobenius notes}  
\author{John E.\ Hutchinson}
\date{\today}


\maketitle

\tableofcontents




%%%%%%%%%%%%%%%%%%%%%%%%%%%%%%%%%%%%
\section{\textbf{Perron Frobenius Theorem}}   I mostly follow \cite{Boy} and~\cite{ Bri}*{pp 57--60}.  See also \cite{Boyd} for some comments and motivation.

%%%%%%%%%%%%%%%%%%%
\subsection{Basic Notions}

\begin{enumerate}

\item  A \emph{primitive} matrix is a square nonnegative matrix some power of which is positive.

\item An \emph{irreducible} matrix is a square nonnegative $ k\times k $ matrix $ A $ such that for every $1\leq  i,j \leq k $ there is a $ p $ such that $ {A^p}_{ij} $ is positive.
\end{enumerate} 

Note that ``primitive $ \Longrightarrow $ irreducible'' but not conversely.

We use the norm $ \|u\|= \sum_i |u_i| $, i.e.\  $ \|u\|= \sum_i u_i $ for nonnegative vectors $ u $, as is always the case here.

%%%%%%%%%%%%%%%%%%%
\subsection{Examples}\mbox{}
\begin{enumerate}
 
\item Let
\[
A=\begin{bmatrix}0&2\\1&1 \end{bmatrix}, \quad B=  \begin{bmatrix}0&4\\1&0 \end{bmatrix},
\]
then $ A^2 $ is positive and so $ A $ is primitive, while $ B $ is irreducible but not primitive.
  

\item \emph{Directed Graphs}.
Suppose $ A $ is a square nonnegative matrix and consider the graph with vertices $ \{1,\dots, k\} $  and a directed edge  $ ij $ iff $ A_{ij}>0 $.  

Then 
\begin{enumerate}

\item $ A $ is primitive iff for some $ p $ there is a path of length $ p $ between every ordered pair of vertices,
\item $ A $ is irreducible iff every ordered pair of vertices can be connected by a path (not all necessarily of the same length).
\end{enumerate} 

\item \emph{Stochastic Matrices}.  Suppose $ A=P $ is a stochastic (or Markov) matrix, i.e.\  $ P $, $ P_{ij}\geq 0 $ and $ \sum_j P_{ij} =1 $.   

The interpretation is that $ P_{ij} $ is the probability of going from node $ i $ to node $ j $ in one time step.  
 
 If $ x^{(0)} $ is the (probability) distribution at time $ 0 $ then $ x^{(0)}P $ is the distribution after one time step   and $ x^{(0)} P^n $ is the  distribution after $ n $  time steps.

\item \emph{Dynamic   Model}.  Here $ A $ determines an evolving mass.

For an  initial mass distribution $  x^{(0)} $ on $ \{1,\dots, k\} $, $ x^{(0)}A $ is the mass distribution after one time step   and 
\[
x^{(n)} :=  x^{(0)} A^n 
\]
 is the   mass distribution after $ n $  time steps.  

The mass   at  node $ i $  is increasing at each time step  if $ \sum_j A_{ij} > 1 $, and is decreasing if  $ \sum_j A_{ij} < 1 $.  

\item \emph{Population and Economic Models}. The mass $ x^{(n)}_i$ in (4) could be interpreted in a \emph{population model} as the number of individuals in group $ i $, and in an \emph{economic model} as the activity level, or resources produced, in sector $ i $.
\end{enumerate} 
 

%%%%%%%%%%%%%%%%%%%%%%%%%
\subsection{Interpretation}  One can probably best  interpret the following theorems in terms of the previous dynamic model.

\medskip If $   x $ is a mass distribution on $ \{1,\dots, k \} $, i.e.\  $ x =[x_1,\dots, x_k] \in \mathbb{R}^k $ with $ x_i \geq 0 $, then $ \|x\| $ is the \emph{total mass} of $x $. 

If $ x $ is the mass distribution   at one time step,  then $ x A =   [(x A)_1,\dots (x A)_k] $ is the mass distribution at the next time step.

Since $ (x A)_j = \sum_ix_i A_{ij} $,  $ A_{ij} $ is interpreted as the proportion or multiple of mass at  $ i $ which is moved to $ j $ at the next time step.  Unlike in the stochastic case we allow $ \sum_j A_{ij} \neq 1 $, i.e.\ we do not assume conservation of mass.  As noted before, the mass   at  node $ i $  is increasing at each time step  if $ \sum_j A_{ij} > 1 $, and is decreasing if  $ \sum_j A_{ij} < 1 $.  

\medskip
In order to interpret  multiplication of $ A $ by a vector on the \emph{right},   consider a \emph{weight function} (value function, utility function, price system)
\[
y = \begin{bmatrix} y_1 \\ \vdots \\ y_k \end{bmatrix} = [y_1,\dots, y_k]^t , 
\]
with values depending on the ``position'' $ i \in \{1,\dots, k\} $.   Then the corresponding \emph{total  weight} (total value, total utility, total price) of the mass distribution $ x $ is $ \langle x,y \rangle = \sum_ix_i y_i $. 

The weight function $ y $ can be quite arbitrary, although normally nonnegative.  For example, consider a markov chain $ P=A $ determining the flow of taxis between three locations in a city.  The weight function for a resident  
living near the first node could possibly look like $ [1, 1/3, 1/4] $.  Multiplying through by a scalar is essentially just a change of units.

\medskip Since 
\[
\langle x, Ay \rangle = \langle x A, y \rangle,\ ( \text{i.e.\ } \langle A^t x, y \rangle ),
\]
one interprets  $ Ay $ as the weight function which acts on an arbitrary mass distribution $ x $ in such a way that the total weight of $ x $ under $ Ay $ is the same as the total weight of  the evolved mass distribution $ xA  $ under $ y $.

Since  $ (Ay)_i = \sum_jA_{ij}y_j $,   $ A_{ij} $ can also be interpreted as  the multiplier of the weight (or price) function at $ j $ which is added to the price at $ i $ to obtain $ (Ay)_i  $.

\medskip
In terms of a directed graph with weights $ A_{ij} $ on the edge \emph{from} $ i $ \emph{to} $ j $, note  that the mass component $ (x A)_j$ is obtained from the multipliers   on the arrows pointing \emph{into} $ j $, whereas the weight (or price) component $ (Ay)_i $ is obtained from the multipliers on the arrows pointing \emph{out of}~$ i $.

\medskip  To interpret the following Perron Frobenius theorem, first note that it applies equally to $ A $ and $ A^t $.

Suppose 
\[
 uA = \lambda u, \quad  Av= \lambda v  .
 \]

Note that for a stochastic matrix, $ \lambda = 1 $, $ u $ is the unique steady state probability (mass) distribution, and $ v = [1,\dots, 1]^t $.  

More generally, $u $ is the unique (up to  scalar multiples) steady state mass distribution modulo rescaling.  The rescaling factor $ \lambda$ is the   total mass multiplier, i.e.\  $ \|uA\| = \lambda \|u\|$.
Similarly, $ v $ is the unique (up to scalar multiples) steady state weight function (price functions) modulo rescaling, and the rescaling is by the same multiplier $ \lambda $.

If $ \beta $  is any other eigenvalue of $ A$ or $ A^t $ then $ | \beta | < \lambda $.   Any right e-vector $ y $ can be written uniquely as  
\[
y =\alpha v + z , \text{ where } z \in \ker u, \text{ i.e.\ } \langle u,z \rangle = 0,
\]
in fact $ \alpha = \langle u,y \rangle / \langle u,v \rangle $.  Moreover, $ \ker u $ is invariant under $ A $, and so the e-values on $ \ker u$ have modulus $ < \lambda $.

It follows that asymptotically
\[ 
A^n \sim \lambda ^n \langle u,v \rangle ^{-1}v \otimes u ,
\]
since $  \langle u,v \rangle ^{-1}v \otimes u $ is just projection onto $ v $ with kernel $ \ker u $.

%%%%%%%%%%%%%%%%%%%%%%%%
\subsection{Main Theorem}
 

 A result analogous to the following is also true in the irreducible case.  See \cite{Boy} and~\cite{ Bri}*{pp 57--60}.  But the main ideas in the proof are already contained in the primitive matrix case.

\begin{theorem}[Perron Frobenius] \label{pft} Suppose $ A $ is primitive.  Then there exists a positive eigenvalue $ \lambda $ such that
\begin{enumerate}
\item $ \lambda  $ is a simple root of the characteristic polynomial %
\footnote{The characteristic polynomial is $ \det (A-xI) $.  The roots $ \lambda $ are the eigenvalues of $ A $.  The algebraic multiplicity of $ \lambda $ equals  the geometric dimension  of the   \emph{generalised} eigenspace corresponding to $ \lambda $, and  so is $ \geq $ the geometric dimension of the eigenspace corresponding to  $\lambda $.}
(so in particular the corresponding eigenspaces for $ A $ and $ A^t $ are one dimensional),
\item $ \lambda  > |\tau| $ if $\tau $ is any other root %
\footnote{\label{arot}Since the coefficients of the characteristic polynomial are real the complex roots come in conjugate pairs.  Suppose $ re^{i\theta} $ is a complex eigenvalue with eigenvector $ w \in \mathbb{C} ^k $.  Let $ w_1 $ and $ w_2 $ be the  real and imaginary parts of $ w $.  Then  taking real and imaginary parts of $ Aw =  re^{i\theta} w $ we  see  in the real 2-dimensional space spanned by $w_1 $ and $ w_2$ that $ A $ is a rotation by $ \theta $ followed by a stretch with $ r $.} 
(so  in particular the spectral radius  of  $ A $ is $ \lambda $),
\item $ A $ has a positive eigenvector $ v $ for $ \lambda $ (and similarly for its transpose $ A^t $), 
\item $ A $ has no nonnegative eigenvectors other than multiples of $v  $  (and similarly for $ A^t $).  
\item   $ \min_i \sum_jA_{ij} \leq \lambda \leq \max_i \sum_j A_{ij} $.
\item Suppose    $ x \geq 0 $.   Then $ Ax < \gamma x \Longrightarrow \lambda < \gamma  $ and $ Ax > \gamma x \Longrightarrow \lambda > \gamma  $.
\end{enumerate} 
\end{theorem} 

\begin{proof}\mbox{}

\noindent \textbf{1}.  \emph{Lemma}: \emph{Suppose $ L: \mathbb{R}^k \to \mathbb{R}^k $ is linear, $ K $ is compact, $ 0 \in K^\circ $ (the interior of $ K $), and $ L^i(K)\subset K^\circ$ for some $ i >0 $.  Then $ \lambda $ is an eigenvalue of $ L $ implies $ | \lambda | < 1 $.}

To see this consider two cases.

If $   | \lambda | > 1 $ then iterates of $ L $ send eigenvectors to $ \infty $, contradicting $   0 \in S^\circ $ and the fact $ L^j(K^\circ) $ is bounded independently of $ j $.

If $ |\lambda| = 1$ consider a corresponding 2-dim subspace $ V $ on which $ L $  acts as a rotation, see footnote~\ref{arot}.  Consider some $ v \in \partial K \cap V $. Either some iterate of $ L $ sends $ v $ to itself, or iterates of $ L $ applied to $ v$ are dense in the unit circle for the subspace $ V $.  Both  contradict  $ L^i(K)\subset K^\circ$.

\medskip \noindent \textbf{2}. Let $\|v\|= \sum_i|v_i| $ for $ v = (v_1,\dots, v_n) \in \mathbb{R}^n$.
  Since $ A $ is nonnegative the map $ f(x) = Ax/\|Ax\| $ sends the unit simplex into itself.  By the Brower fixed point theorem there is a fixed point $ v $, and $ v $ is  then a  nonnegative eigenvector of $ A $ with   eigenvalue $ \lambda >0$.  Since some iterate of $ A $ is positive, $v  $ is positive.

This proves (3).


\medskip \noindent \textbf{3}.  Use $ v $ and $ \lambda $ to transform $ A $ to a  \emph{stochastic}  (i.e.\ probability)  matrix $ P $ with right eigenvector $(1,\dots, 1)^T $ and eigenvalue $ 1 $.   That is
\begin{equation} \label{dfPA}
P := \lambda ^{-1} V^{-1} A V, \text{ or equivalently } P_{ij} = \frac{A_{ij}v_j}{\lambda v_i} = \lambda ^{-1}\frac{v_j}{v_i} A_{ij},
\end{equation} 
where $ V $ is the diagonal matrix with entries from $ v $.   That is, rescale the axes to obtain a coordinate change so that   $ v $ becomes  the unit vector, and then rescale $ A $ so the corresponding  eigenvalue is  one.

Note that 
\begin{align*} 
0\leq P_{ij}, \quad \sum_j P_{ij} = 1,\\
\lambda ^{-n}  \det(A-\tau I) &= \det( P-\lambda ^{-1}\tau I) \\
Av = \tau v \Longleftrightarrow  P(V^{-1}v) &= \lambda ^{-1}\tau\, V^{-1}v 
\end{align*} 

\medskip \noindent \textbf{4}. It is now sufficient to prove (1) and (2) for $ P $, with $ \lambda   $ replaced by $ 1 $ and corresponding  e-vector $ (1,\dots,1)^T  $.  In fact we will prove  (1) and (2)  for   the transpose $ P^t $, i.e.\ for the left action of $ P $. 

First note that the \emph{transpose} $ P^t $ maps the unit simplex $S $  into itself.  Let $ w \in S $ be the unique nonnegative left e-vector with eigenvalue one for $ P^t $.%
\footnote{This eigenvector exists either from \textbf{1} applied to the transpose of $ P $, or by the standard markov chain result for the existence of a stationary probability distribution.}


 Consider the translate $ S - w $. Then $ P^t :S-w\to S-w $.
 
 Working in the codimension one subspace spanned by $ S-w $ we can apply step \textbf{1}   to $ P^t $, noting that since for some power of $ P^t $ all entries are positive,  this power maps $ S-w $ to its interior.  Hence all eigenvalues of $  P^t $ restricted to this subspace are strictly less than 1.  The remaining e-value and e-vector of $ P^t $ are $ 1 $ and $ w $. 
 
 This gives (1) and (2) for $ P^t  $, and hence for $ P $, and hence   for $ A $.
 
 \medskip \noindent \textbf{5}.   To prove  (4) suppose $ w  $ is a nonnegative e-vector of $ A $ with e-value $  \beta$.  Since $ A^i w = \beta ^i w $ and the left side is positive for some $ i $, it follows $ w $ is in fact positive and so $ \beta $ is real and positive.  Also, $ \beta \leq \lambda $ since $ \lambda $ is the spectral radius.
 
 Suppose $ v $ is a positive e-vector corresponding to the   e-value $ \lambda $ and choose a positive e-vector $ w $ for $ \beta $ such that $ w>v $ componentwise.   
 Then $ A^iw >  A^iv $ for all $ i $, hence $ \beta ^i w > \lambda ^i v $ for all $ i $, which is impossible if $ \beta < \lambda $.
 
Hence $ \beta = \lambda $.  Since  $ \lambda $ is a simple root of the characteristic equation the corresponding eigenspace is one dimensional and so we are done.

 \medskip \noindent \textbf{6}.  For (5) let $ v_p = \min\{v_1,\dots v_k\}$. Then
 \[
 \lambda v_p = \sum_j A_{pj} v_j \geq \bigg( \sum_j A_{pj} \bigg) v_p \geq \bigg( \min_i \sum_p A_{ij}\bigg) v_p.
 \]
 Similarly for $ \max $.
 
 \medskip \noindent \textbf{7}.    Suppose $ Ax < \gamma x $ for some $ x \geq 0 $. It follows $ \gamma > 0 $ and so $ x > 0 $.
 
 W.l.o.g.\ suppose $ v<x $ where $ Av = \lambda v $.  Note $ v>0$.
 
 Then for all $ n $, 
 \[
 \lambda ^n v = A^n v < A^n x < \gamma ^n x.
 \]
 This is impossible if $ \gamma < \lambda $, hence $ \lambda \leq \gamma $.
 
 But $ Ax < \gamma x $ and $ x>0 $ implies $ Ax < (\gamma- \epsilon) x $ for some $ \epsilon > 0 $.
 Hence $ \lambda \leq \gamma - \epsilon  $ and so $ \lambda < \gamma $.   
 
  A similar proof applies if $ Ax > \gamma x $ for some $ x \geq 0 $.
 \end{proof} 


\begin{remark}
***to be checked***  The fact that $ P $, defined by ``normalising'' $ A $  by its right eigenvector as in \eqref{dfPA}, is a stochastic matrix, is 
important in understanding the Ruelle Perron Frobenius theorem, particularly the existence of a \emph{shift invariant} equilibrium measure.  

*** (Need normalising here) *** It is easily checked that the left e-vector $ \pi $ of $  P $ is the same as the left e-vector $ u $ of $ A $ \emph{weighted} by its right e-vector $ v $, which is the same as the diagonal vector of $v\otimes u/\langle u,v \rangle  $, see~\eqref{limA}.   That is,
\[
\pi = (u_1v_1,\dots, u_kv_k) .
\]
Of course, the right e-vector if $ P $ is just $ (1,\dots, 1) $.
\hfill $\square$  \end{remark} 


%%%%%%%%%%%%%%%%%%%%%%%%%%%%
\subsection{Limit Theorem}
\begin{corollary}\label{corlt}
  Suppose $ A $ is primitive, $  \lambda  $ is the spectral radius, $ u  $ is a
 left (row) e-vector for $ \lambda   $, and $ v $ is a right (column) e-vector for $ \lambda   $.

Then 
\begin{equation} \label{limA} 
  \lambda ^{-n}A^n \to \frac{v\, u }{u  v}  = \frac{v \otimes u}{  \langle u , v \rangle }  \text{ exponentially fast as } n\to \infty ,
\end{equation} 
  where $ v  u  = v\otimes u $ is the rank one $ k \times k $ matrix obtained by multiplying the column vector  $v  $ by the row vector~$ u$, and $ uv = \langle u , v \rangle $ is the scalar obtained from matrix multiplication in the reverse direction.
\end{corollary}


\begin{remark}\mbox{}

\begin{enumerate}
\item The matrix $ \dfrac{v\, u}{uv} $ is projection onto the subspace spanned by $ v $ whose kernel is the codimension one subspace annihilated by $ u $.

\item W.l.o.g.\  we \emph{normalise} the e-vectors $ u $ and $ v $ so 
\[
\|u\|:= \sum_i|u_i| = 1 , \quad  u v = \sum_i u_i v_i =1 , 
 \]
 \item 
In particular,  $  \lambda ^{-n}A ^n   \to  v \, u $ and  $  \lambda ^{-n}A ^n_{ij}  \to  v_i\, u_j $, exponentially fast. 


\item  Taking norms, $   \lambda  ^{-n} \|A\|^n \to \|v\, u\| $.  Taking logs and dividing by $ n $, 
\[   \lim_{n\to \infty} \frac{1}{n} \log \|A^n\| =  \lambda . \]

\item It follows from (3) that
\begin{align*}
\lambda ^{-n} x^{(n)} &\to \langle x^{(0)},v\rangle\, u ,  \quad \text{since } x^{(n)} = x^{(0)} A^n, \\
\lambda ^{-n} \|x^{(n)}\| &\to \langle x^{(0)},v\rangle,  \quad \text{since }  \|u\|=1,  \\
\frac{x^{(n)}}{\|x^{(n)}\|}  &\to u, \quad \text{from the preceding two limits},     \\
\frac{x^{(n+1)}_i}{ x^{(n)}_i}  &\to \lambda,  \quad \text{from the first limit} . 
\end{align*} 

\item The first limit above says the asymptotic distribution is  $ u $ scaled by $ \langle x^{(0)},v\rangle  $.
The second says the asymptotic total mass is $ \langle x^{(0)},v\rangle $.

If $ v $ is the vector of units as in the stochastic case then $ \langle x^{(0)},v\rangle = \sum_i x^{(0)}_ i  $.
In general the \emph{interpretation of $ v $} is that it gives the contribution of each component of $  x^{(0)}$ to the asymptotic mass distribution.   

The \emph{interpretation of $ u $} by the third limit is as  the limit of the normalised distributions.  For the Markov case it is the steady state probability distribution.

The fourth limit says that the one time step limit growth factor in every component is $ \lambda $.
\end{enumerate} 
\end{remark}

\begin{proof}[Proof of Corollary] W.l.o.g.\ take $ \lambda=1 $ by multiplying $ A $ by a constant.  Then all other e-values have modulus less than $ 1 $.

Choose a basis such that  $ A $ is in Jordan normal form with  $ 1 $ the top left entry and all other diagonal entries strictly less than $ 1 $ in modulus.  

Since $ 1 $ is a simple root of the characteristic equation, all entries in  the left column and top row other than the unit diagonal element must equal zero.  It follows by uniqueness of the left and right eigenspaces corresponding to $1  $ that in the new basis $ u= (1,0,\dots, 0) $ and $ v= (1,0,\dots, 0)^t $.

Moreover, since all diagonal elements are less than one in modulus,  it is easy to check that powers of any submatrix corresponding to any other generalised e-space converges exponentially fast to $ 0 $.  For example, 
\[
B = \begin{bmatrix} \beta &1 &0 \\ 0 &\beta & 1 \\ 0 & 0 & \beta \end{bmatrix}, \quad
B^2 = \begin{bmatrix} \beta^2 &2 \beta  &1 \\ 0 &\beta^2  & 2 \beta  \\ 0 & 0 & \beta^2 \end{bmatrix}, \quad
B^3 = \begin{bmatrix} \beta^3 &3 \beta^2  &3 \beta  \\ 0 &\beta^3  & 3 \beta^2  \\ 0 & 0 & \beta^3 \end{bmatrix}, 
\quad \text{etc.}
\]

Hence $ A^n $ converges exponentially fast to the matrix with $ 0 $ everywhere except for $ 1 $ in the top left corner.  Thus for  this coordinate system, $ A^n \to v\,   u^t $.   Since the assertion is independent of the coordinate system,  we are done.
\end{proof} 


%%%%%%%%%%%%%%%%%%%%%%%%%%%%%%%%
%%%%%%%%%%%%%%%%%%%%%%%%%%%%%%%%
%%%%%%%%%%%%%%%%%%%%%%%%%%%%%%%%
%%%%%%%%%%%%%%%%%%%%%%%%%%%%%%%%


\begin{thebibliography}{9} 

%\bib{Bed}{article}{
%   author={Bedford, Tim},
%   title={Applications of dynamical systems theory to fractals---a study of
%   cookie-cutter Cantor sets},
%   conference={
%      title={Fractal geometry and analysis},
%      address={Montreal, PQ},
%      date={1989},
%   },
%   book={
%      series={NATO Adv. Sci. Inst. Ser. C Math. Phys. Sci.},
%      volume={346},
%      publisher={Kluwer Acad. Publ.},
%%      place={Dordrecht},
%   },
%   date={1991},
%   pages={1--44},
%%   review={\MR{1140719}},
%}

%\bib{Bil}{book}{
%   author={Billingsley, Patrick},
%   title={Ergodic theory and information},
%   publisher={John Wiley \& Sons Inc.},
%%   place={New York},
%   date={1965},
%   pages={xiii+195},
%%   review={\MR{0192027 (33 \#254)}},
%}

\bib{Boyd}{article}{
   author={Boyd, Stephen},
   title={Perron Frobenius Theory},
  note={Online Notes, Chapter 17 of course notes} 
     }


\bib{Boy}{article}{
   author={Boyle, Mike},
   title={Notes on the Perron-Frobenius Theory of Nonnegative matrices},
  note={Online Notes} 
     }


\bib{Bri}{book}{
   author={Brin, Michael},
   author={Stuck, Garrett},
   title={Introduction to dynamical systems},
   publisher={Cambridge University Press},
   place={Cambridge},
   date={2002},
   pages={xii+240},
%   isbn={0-521-80841-3},
%   review={\MR{1963683 (2003m:37001)}},
}
	
%\bib{Edg}{book}{
%   author={Edgar, Gerald A.},
%   title={Measure, topology, and fractal geometry},
%   series={Undergraduate Texts in Mathematics},
%   publisher={Springer-Verlag},
%   place={New York},
%   date={1990},
%   pages={xiv+230},
%%   isbn={0-387-97272-2},
%%   review={\MR{1065392 (92a:54001)}},
%}

%\bib{Fal}{book}{
%   author={Falconer, Kenneth},
%   title={Techniques in fractal geometry},
%   publisher={John Wiley \& Sons Ltd.},
%%   place={Chichester},
%   date={1997},
%   pages={xviii+256},
%%   isbn={0-471-95724-0},
%%   review={\MR{1449135 (99f:28013)}},
%}

%\bib{Kea}{article}{
%   author={Keane, Michael S.},
%   title={Ergodic theory and subshifts of finite type},
%   conference={
%      title={Ergodic theory, symbolic dynamics, and hyperbolic spaces
%      (Trieste, 1989)},
%   },
%   book={
%      series={Oxford Sci. Publ.},
%      publisher={Oxford Univ. Press},
%%      place={New York},
%   },
%   date={1991},
%   pages={35--70},
%%   review={\MR{1130172}},
%}

%\bib{Kea72}{article}{
%   author={Keane, Michael},
%   title={Strongly mixing $g$-measures},
%   journal={Invent. Math.},
%   volume={16},
%   date={1972},
%   pages={309--324},
%%   issn={0020-9910},
%%   review={\MR{0310193 (46 \#9295)}},
%}

%\bib{MW}{article}{
%   author={Mauldin, R. Daniel},
%   author={Williams, S. C.},
%   title={Hausdorff dimension in graph directed constructions},
%   journal={Trans. Amer. Math. Soc.},
%   volume={309},
%   date={1988},
%%   number={2},
%   pages={811--829},
%%   issn={0002-9947},
%%   review={\MR{961615 (89i:28003)}},
%}

%
%\bib{Mil}{article}{
%   author={Milnor, John},
%   title={Dynamics: Introductory Lectures},
%  note={Preliminary Notes} 
%     date={2001}
%     }

%%\bib{Mis1}{article}{
%%   author={Misiurewicz, Micha{\l}},
%%   title={A short proof of the variational principle for a ${\bf
%%   Z}\sb{+}\sp{N}$ action on a compact space},
%%   conference={
%%      title={},
%%      address={Rennes},
%%      date={1975},
%%   },
%%   book={
%%      publisher={Soc. Math. France},
%% %     place={Paris},
%%   },
%%   date={1976},
%%   pages={147--157. Ast\'erisque, No. 40},
%%%   review={\MR{0444904 (56 \#3250)}},
%%}
%	

%\bib{Mis1}{article}{
%   author={Misiurewicz, M.},
%   title={A short proof of the variational principle for a $Z\sp{n}\sb{+}$\
%   action on a compact space},
% %  language={English, with Russian summary},
%   journal={Bull. Acad. Polon. Sci. S\'er. Sci. Math. Astronom. Phys.},
%   volume={24},
%   date={1976},
%%   number={12},
%   pages={1069--1075},
%%   issn={0001-4117},
% %  review={\MR{0430213 (55 \#3220)}},
%}
%		

%

%\bib{PY}{book}{
%   author={Pollicott, Mark},
%   author={Yuri, Michiko},
%   title={Dynamical systems and ergodic theory},
%   series={London Mathematical Society Student Texts},
%   volume={40},
%   publisher={Cambridge University Press},
%%   place={Cambridge},
%   date={1998},
%   pages={xiv+179},
%%   isbn={0-521-57294-0},
%%   isbn={0-521-57599-0},
%%   review={\MR{1627681 (99f:58130)}},
%}
%	

%

%
%\bib{Walop}{article}{
%   author={Walters, Peter},
%   title={Ruelle's operator theorem and $g$-measures},
%   journal={Trans. Amer. Math. Soc.},
%   volume={214},
%   date={1975},
%   pages={375--387},
%%   issn={0002-9947},
%%   review={\MR{0412389 (54 \#515)}},
%}
%		

%\bib{Walvp}{article}{
%   author={Walters, Peter},
%   title={A variational principle for the pressure of continuous
%   transformations},
%   journal={Amer. J. Math.},
%   volume={97},
%   date={1975},
%   number={4},
%   pages={937--971},
%%   issn={0002-9327},
%%   review={\MR{0390180 (52 \#11006)}},
%}
%		

%\bib{Wal}{book}{
%   author={Walters, Peter},
%   title={An introduction to ergodic theory},
%   series={Graduate Texts in Mathematics},
%   volume={79},
%   publisher={Springer-Verlag},
%%   place={New York},
%   date={1982},
%   pages={ix+250},
%%   isbn={0-387-90599-5},
%%   review={\MR{648108 (84e:28017)}},
%}

%\bib{Wil}{article}{
%   author={Williams, Susan G.},
%   title={Introduction to symbolic dynamics},
%   conference={
%      title={Symbolic dynamics and its applications},
%   },
%   book={
%      series={Proc. Sympos. Appl. Math.},
%      volume={60},
%      publisher={Amer. Math. Soc.},
%%      place={Providence, RI},
%   },
%   date={2004},
%   pages={1--11},
%%   review={\MR{2078843 (2005e:37024)}},
%}
%		
%\bib{You}{article}{
%   author={Young, Lai-Sang},
%   title={Entropy in dynamical systems},
%   conference={
%      title={Entropy},
%   },
%   book={
%      series={Princeton Ser. Appl. Math.},
%      publisher={Princeton Univ. Press},
%      place={Princeton, NJ},
%   },
%   date={2003},
%   pages={313--327},
%%   review={\MR{2035829 (2004k:37010)}},
%}

%	

\end{thebibliography}

\end{document} 

